\subsection{* 有界连续函数空间的情形}

对每个拓扑空间 X，设 \(\mathcal{C}^{b}(X)\) 表示从 X 到 K 中的所有连续且有界的映射所成的 Banach 空间，赋以由下式定义的范数

\[
\| f \| = \sup _ {x \in X} | f (x) |\tag{6}
\]

(GT, X, § 3, No.~2)。当 X 紧时，X 上的每个连续函数都有界 (GT, IV, § 6, No.~1)，我们把 \(\mathcal{C}(X)\) 记作 \(\mathcal{C}^b(X)\)。

在本节及下一节中，我们将用到下面的引理，它是 Lebesgue 定理 (INT, IV, 第 2 版, § 4, No.~3, 定理 2) 的一个特例，这里用到把 \(\mathcal{C}(\mathrm{X})'\) 的元素解释为 X 上的测度。

引理 2. --- 设 X 是一个紧空间。若序列 \((f_n)_{n\in \mathbf{N}}\) 在 \(\mathcal{C}(\mathbf{X})\) 中有界且在 X 上简单收敛到连续函数 \(f\)，则 \(\mu(f) = \lim_{n\to \infty}\mu(f_n)\) 对每个 \(\mu\)（属于 \(\mathcal{C}(\mathbf{X})'\)）成立。

命题 3. --- 设 X 是一个紧空间，A 是 \(\mathcal{C}(\mathrm{X})\) 的一个有界子集。为使 A 对于简单收敛拓扑相对紧，必须且只需它对于 \(\sigma(\mathcal{C}(\mathrm{X}), \mathcal{C}(\mathrm{X})')\) 相对紧。

简单收敛拓扑是 Hausdorff 的且比 \(\sigma(\mathcal{C}(X), \mathcal{C}(X)')\) 粗，故所述条件是充分的 (GT, I, § 9, No.~4, 推论 3)。

现在假设 A 对于简单收敛拓扑相对紧。设 \((f_{n})_{n\in\mathbf{N}}\) 是 A 的元素的一个序列。由命题 2 (IV, 第 33 页)，存在一个序列 \((f_{n_{k}})\)（从 \((f_{n})\) 中抽取），它简单收敛到一个连续函数 f。由引理 2，有界序列 \((f_{n_{k}})\) 对于 \(\sigma(\mathcal{C}(\mathrm{X}),\mathcal{C}(\mathrm{X})')\) 趋于 f。于是 Šmulian 定理 (IV, 第 36 页, 定理 2) 表明 A 对于 \(\sigma(\mathcal{C}(\mathrm{X}),\mathcal{C}(\mathrm{X})')\) 相对紧。

推论. --- 设 S 是一个拓扑空间，A 是 \(\mathcal{C}^{b}(S)\) 的一个有界子集。下列条件等价：

\begin{enumerate}
\def\labelenumi{(\roman{enumi})}
\item
  A 对于 \(\sigma(\mathcal{C}^b(\mathbf{S}),\mathcal{C}^b(\mathbf{S})')\) 相对紧；
\item
  若 \((f_n)_{n\in \mathbb{N}}\) 是 \(\mathbf{A}\) 的元素的序列，\((x_m)_{m\in \mathbb{N}}\) 是 \(\mathbf{S}\) 的点的序列，且累次极限
\end{enumerate}

\[
\gamma = \lim _ {m \rightarrow \infty} \lim _ {n \rightarrow \infty} f _ {n} (x _ {m}), \quad \delta = \lim _ {n \rightarrow \infty} \lim _ {m \rightarrow \infty} f _ {n} (x _ {m})
\]

存在，则 \(\gamma = \delta\)。

设 X 是 S 的 Stone-Čech 紧化 (GT, IX, § 1, No.~6)，\(\alpha\) 是从 S 到 X 中的典范映射。令 \(D = \alpha(S)\)。映射 \(\phi: f \mapsto f \circ \alpha\) 是从赋范空间 \(\mathcal{C}(X)\) 到赋范空间 \(\mathcal{C}^b(S)\) 上的一个同构；令 \(\tilde{A} = \phi^{-1}(A)\)。由于 X 紧且 D 在 X 中稠密，命题 2 (IV, 第 33 页) 表明条件 (ii) 等价于 \(\tilde{A}\) 对于简单收敛拓扑的紧性。(i) 与 (ii) 的等价性于是由命题 3 得出。*

